\section{Short blocks, rational shifts, and the core estimate}
\label{sec:arithmetic}

Starting from the localization proved in the preceding section, we now
split the arguments into side terms and core terms.  Side terms are controlled by
square estimates for projection increments, together with a comparison on
a bounded block of depths.  Core terms are controlled by separation of
the frequencies left after removing a common quadratic phase.  We prove
the coefficient-dependent parts of both arguments here.  The projection,
graph, and amplitude summations that we use from \cite{H,V} are identified
at their points of application.

Throughout this section, the primitive integer vector
\(b=(0,b_1,b_2,b_3)\), the integers \(e_j\), and the nonzero integers
\(c_j\) are those of Lemma~\ref{lem:algebra}.  Constants may depend on
these fixed integers and on the admissible kernel bounds.  We write
\(e(t)=\exp(2\pi i t)\).  For an interval \(I\), the notation
\(\|f\|_{2,I}\) denotes \((|I|^{-1}\int_I|f|^2)^{1/2}\), also for
Hilbert-space-valued functions.

\subsection{A comparison on a short block}

We first establish the version of \cite[Lemma~7.2]{H} needed by the
side argument.  This estimate compares randomized inputs with the
fixed-input constant \(C_N(q)\); in particular, the exponent of the
block length must tend to zero as \(q\) tends to two.

\begin{proposition}[Short-block comparison]\label{prop:short-block-b}
Let \(2<q<3\).  Let a finite dyadic tree piece have root \(R\), with
internal summation intervals in at most \(K\) consecutive depths and
within the original range of at most \(N\) depths.  Its terminal
intervals, called leaves, partition \(R\) and occur in this depth block
or one depth below it.  In role \(j\), let \(\mathbf z_j\) be a
deterministic Hilbert-space-valued function supported on \(R\), with
\(\|\mathbf z_j\|_{2,L}\le s_j\) on every leaf \(L\).
Suppose scalar randomizations \(z_j^\sigma\) satisfy
\[
 \bigl(\mathbb E_\sigma|z_j^\sigma(y)|^a\bigr)^{1/a}
 \le C_a\|\mathbf z_j(y)\|
\]
for \(a=2,4,8\), with constants independent of the tree and \(K\).
For measurable signs
\(|\epsilon_I(\sigma)|\le1\) and any subcollection of internal intervals,
\[
 \mathbb E_\sigma\left|
 \sum_I |I|\epsilon_I(\sigma)H_I(z_0^\sigma,z_1^\sigma,
                                  z_2^\sigma,z_3^\sigma)\right|
 \le (1+C_N(q))C_{q,b}(1+\log(K+2))^C
       (K+2)^{11(q-2)/(q-1)}|R|\prod_j s_j .
\]
The constant \(C_{q,b}\) may also depend on the fixed moment
bounds \(C_a\) and the kernel bounds.
\end{proposition}

\begin{proof}
Homogeneity reduces to \(s_j=1\); a zero size makes the form zero.
Multiplying all signs by one random scalar of modulus one reduces the
expected absolute value to the absolute value of an expected signed sum.

We use the partition of the four leaf depths from the proof of
\cite[Lemma~7.2]{H}.  Here is the information from that partition that
will be needed.  If the least depth occurs more than once, designate two
roles attaining it.  If it occurs once, designate that role and
separate the least and second least depths by the smallest dyadic
interval of depth indices containing them in opposite halves.
Specify the nonempty set of roles whose depths belong to its upper
half, and designate one of those roles by a fixed ordering.
Every role that is neither the minimum role nor in that specified set
has depth greater than the upper endpoint of that dyadic interval.
This partitions the depth tuples into rectangles of masks
\(E_j\), each a union of leaves.  There are
\(O(1+\log(K+2))\) length classes and a bounded number of role assignments.
Within each class, there are \(O(K+1)\) positions; the masks in each of
the two designated roles, say \(a,a'\), are disjoint over those positions.
For a fixed rectangle there is an index \(t\) such that every role has
leaf depth at least \(t\), with the possible exception of role \(a\).
Only summation depths \(s<t\) contribute.  Indeed a leaf meeting an
internal interval must be its strict descendant.  The same assertion
holds after the small coordinate enlargements below, by the kernel's
fixed interior support margin.

Write \(r_s\) for the length at depth \(s\), and set
\(\delta=\beta r_t\), where \(\beta>0\) is a sufficiently small fixed
dyadic number depending on \(b\) and the support margin.  Decompose the
physical coordinates as
\[
 x=\delta(n+u),\qquad \tau=\delta(m+v),\qquad
 n_j=n+b_jm,\qquad w_j=u+b_jv.
\]
The slow coordinates \(n,m\) are integers.  Choose the fast
representatives from the parallelogram
\begin{equation}\label{eq:fast-tile}
 \mathcal P_b=
 \{(u,v):0\le u<1,\quad 0\le u+b_1v<|b_1|\}.
\end{equation}
This is a fundamental domain for translation by \(\mathbb Z^2\), of
area one.  To decompose a point, first choose \(n\) so that its first
coordinate minus \(n\) belongs to \([0,1)\), and then choose \(m\) so
that \(u+b_1v\) belongs to \([0,|b_1|)\).  The latter choices differ
by \(b_1\), so existence and uniqueness hold for either sign of \(b_1\).

Let \(\mathbb T_0=\mathbb R/(T_0\mathbb Z)\), with Haar probability
measure, where \(T_0\) will be a sufficiently large fixed number.
Define \(F_j^\sigma(p,w)\) by sampling
\(z_j^\sigma(\delta(p+w))\mathbf1_{E_j}(\delta(p+w))\)
on a bounded interval of real lifts and setting it to zero elsewhere
on \(\mathbb T_0\).  Use lift intervals \([0,1)\) and
\([0,|b_1|)\) in roles zero and one, and bounded intervals containing
the ranges of \(w_2,w_3\) on \(\mathcal P_b\) in the other roles.

The fast Haar integral reproduces integration over \(\mathcal P_b\),
up to its fixed normalization.  To check the potentially noninjective
pair coordinates, suppose a torus progression has all four bounded
lifts \(\widetilde w_j\) in these windows.  Put
\[
 \widetilde u=\widetilde w_0,\qquad
 \widetilde v=\sum_j e_j\widetilde w_j.
\]
The integer inverse of Lemma~\ref{lem:algebra} gives
\(\widetilde w_j=\widetilde u+b_j\widetilde v\pmod {T_0}\).
All the discrepancies are bounded in terms of the fixed windows and
\(b,e\).  Choose \(T_0\) larger than these bounds.  Each discrepancy
then vanishes, so these are equalities of real numbers and
\((\widetilde u,\widetilde v)\in\mathcal P_b\).  Conversely, every
point of the tile gives exactly these lifts.  This verifies the sampling
identity without asserting that any coordinate pair is unimodular.

Replace the kernel by its value at \((\delta n,\delta m)\).
At depth \(s\), the error is \(O_b(\delta/r_s)\), with a fixed support
enlargement.  The normalized sampled form has factor
\((\delta/r_s)^2\).  On the product of local counting measure and
the fast torus, any two distinct coordinate maps are injective on the
lattice and Haar-surjective on the torus.  Consequently the two-
\(L^1\), two-\(L^\infty\) bounds hold, and interpolation gives the
absolute form estimate whenever the four reciprocal exponents sum to
at most two.

For clarity, let \(M_{j,I}\) be the squared mass of \(\mathbf z_j\)
in all masked sampling windows relevant to \(I\), including their
bounded overlap multiplicity.  Choosing \(\beta\) small enough puts
these windows inside \(I\).  The leaf normalization therefore gives
\[
 M_{j,I}\le C_b|I|,\qquad
 \sum_{\operatorname{depth}(I)=s}M_{j,I}\le C_b|E_j|.
\]
The absolute \(L^2\) estimate and the moment hypothesis bound the
sampling errors at depth \(s\) by
\[
 C_b\frac{\delta}{r_s}|E_a|^{1/2}|E_{a'}|^{1/2}.
\]
The sum over \(s<t\) is bounded, and summing rectangle positions uses
the disjointness of both designated masks.

We next improve the integrability of the fast inputs.  On their fast
Fourier coefficients define
\[
 \widehat{U_{j,\gamma}F}(p,k)
   =e(\gamma k^2/c_j)\widehat F(p,k),
 \qquad \gamma\in\mathbb R/(C_c\mathbb Z),
 \quad C_c=\operatorname{lcm}_j|c_j|.
\]
The fast progression integral is unchanged under simultaneous
application of these unitaries.  Its Fourier support obeys
\(\sum k_j=\sum b_jk_j=0\), so Lemma~\ref{lem:algebra} gives
\(k_j=c_j(h+b_jh')\) for real \(h,h'\).  Thus
\[
 \sum_j k_j^2/c_j=\sum_jc_j(h+b_jh')^2=0.
\]
For a single role, integrating the fourth power first in the fast
variable and then in \(\gamma\) equates the sums and the sums of
squares of two pairs of Fourier indices.  Those pairs agree as
multisets, whence
\begin{equation}\label{eq:fast-fourth}
 \mathbb E_\gamma\int_{\mathbb T_0}|U_{j,\gamma}F(p,w)|^4\,dw
 \le2\|F(p,\cdot)\|_2^4.
\end{equation}
Fourier polynomial approximation is justified by the pairwise
\(L^2\) form estimate and this same fourth-moment estimate for differences.

Let \(\mathfrak e_j(p)\) now denote the fast \(L^2\) size of the deterministic
stack in the sampling window.  Window overlap gives
\(\delta\sum_p \mathfrak e_j(p)^2\le C_b|E_j|\).
If all leaves of a role have depth at least \(t\), then
\(\mathfrak e_j(p)\le C_b\): disjoint leaves meeting a window of length
\(O_b(\delta)\) have total length at most \(O_b(\delta)+2r_t\).
In every role, including the possible remaining role \(a\),
\[
 \frac{\delta}{|I|}\sum_{p\ \text{relevant to}\ I}\mathfrak e_j(p)^2\le C_b.
\]
Minkowski's inequality transfers the assumed sign moments to these
window sizes.  Intermediate moment orders below eight are bounded
by the eighth moment using H\"older's inequality.

Write \(G_j=U_{j,\gamma}F_j\).  If role \(a\) has leaves above
depth \(t\), clip it radially at \(A=(K+2)^{10}\).  The clipped tail
has local \(L^1\) norm at most \(A^{-1}\|F_a\|_2^2\), whose
\(L^4\) norm in the random parameters is \(O_b(A^{-1})\).
In the other roles the local
\(L^3\) norms have bounded fourth moments, by
\eqref{eq:fast-fourth} and the preceding window bounds.  The absolute
form estimate with exponents \((1,3,3,3)\) therefore gives expected
tail error \(O_b(A^{-1})\) per normalized cell.  Even summing
\(O((K+1)^2)\) positions and depths is harmless.
For the rest of the proof, replace \(G_a\) by this clipped function
when that role is present.

To return to the original continuous form, fix the random parameters
and \((u,v)\in\mathbb T_0^2\), and plant bursts
\[
 g_j(y)=\sum_pG_j(p,u+b_jv)
                   \mathbf1_{\{|y-\delta p|<c'\delta\}}.
\]
Choose \(c'>0\) so small that
\(c'(2+\max_j|b_j|\sum_l|e_l|)<1\).
If a real configuration lies in bursts centered at \(p_j\), set
\(n=p_0\), \(m=\sum_j e_jp_j\).  Then
\[
 |x/\delta-n|<c',\qquad
 |\tau/\delta-m|<c'\sum_j|e_j|,
 \qquad |p_j-n-b_jm|<1.
\]
The last quantity is an integer and hence is zero.  For this slow
pair the allowed offsets form the polygon
\(\{|u'+b_jv'|<c'\ (0\le j\le3)\}\), of fixed positive finite area.
The center-valued kernel integral on the polygon therefore gives
exactly the sampled product, times that area and \(\delta^2\).
The second kernel replacement has the same error estimate as the
first: unitarity and clipping preserve fast \(L^2\) bounds and do
not enlarge the set of slow indices.  We can now apply \(C_N(q)\)
to these fixed compactly supported burst inputs, using the original
admissible kernels.

For completeness, the exponent calculation is unchanged.  Put
\[
 p_a=p_{a'}=p_*=2+2(q-2),\qquad
 p_j=\frac{2p_*}{p_*-2}\quad(j\notin\{a,a'\}).
\]
Use fast exponent \(p_*\) in the clipped role and exponent four
in the other roles.  Fast pairwise H\"older, random H\"older of
exponent four, spatial H\"older with the \(p_j\), and the mixed-norm
maximal theorem \cite[Lemma~7.1]{H} apply: all exponents in that
maximal theorem exceed \(q\).  The resulting regular-role factors
are \(C_b|E_j|^{1/p_j}\); the clipped factor is
\(C_{q,b}A^{1-2/p_*}|E_a|^{1/p_*}\).
Summing \(O(K+1)\) positions in the two disjoint designated masks
costs \((K+2)^{1-2/p_*}\).  Hence the total power is
\[
 A^{1-2/p_*}(K+2)^{1-2/p_*}
       =(K+2)^{11(q-2)/(q-1)}.
\]
The logarithmic number of length classes finishes the proof.
\end{proof}

This proposition supplies exactly the input used by the side estimates
\cite[Lemmas~7.3--7.4 and Proposition~7.5]{H} and
\cite[Section~6.5]{V}.  Their remaining arguments concern nested
projections, frozen stacks, and variation in depth.  In addition to the
local absolute form bound, they assume the square estimate H(17) for
the linear sides.  We identify its source below; no coordinate
substitution is needed in those projection arguments.

\subsection{A common gauge and a lattice of rational shifts}

We describe the data from the mode construction of
\cite[Section~8]{H} before checking its arithmetic.  Fix a localized
group \(G\), with central tag \(z=(\theta_z,C_z,\lambda_z)\), and
a role \(l\).  A \emph{primary} is an available main column of this
role's outer projection.  Its first matching cell is the cell where
the finite expansion described below is chosen; that choice and its
labels are then inherited on descendants.  Thus later freezing does
not resample a primary's expansion.

Choose an integer \(Q_0\) clearing the denominators of \(C_z\), and
put \(\vartheta=\theta_z/Q_0\).  For an offset \(\rho\) in the unit
cube of horizontal coordinates, define
\[
 d^\rho(y)=\rho+\{\theta_z y-\rho\},\qquad
 p^\rho(y)=\lambda_z y^2+(\theta_z y)^TC_zd^\rho(y).
\]
Choose \(\delta_c=\exp(-k^{C_{\rm cut}})\) sufficiently small after
the main budgets, with fixed exponent \(C_{\rm cut}\).
Let \(\zeta^\rho\) be the product Gevrey cutoff of
\cite[Section~8]{H}, vanishing within distance \(\delta_c\) of each
branch cut and equal to one outside distance \(2\delta_c\).
It is bounded between zero and one.  For each role and group,
one offset is kept fixed throughout its history.  Its expectation
is at least \(1/2\), and
\(\mathbb E_\rho|1-(\zeta^\rho)^2|^2\le C d_{\max}\delta_c\),
where \(d_{\max}\) denotes an upper bound for the dimensions.

For a primary \(\Phi=(\varphi_\Phi,\epsilon_\Phi^{1/2}\mathbf e_\Phi)\)
with normalized tag \(v\),
the endpoint splitting from the localization section and the chart
expansion of \cite[Lemma~4.4]{H} give, at its first match,
\[
 \zeta^\rho\varphi_\Phi e(-c_l p^\rho)
   =\sum_m a_{\Phi,m}e(\lambda_my^2+\gamma_my)+O(\delta_c).
\]
The chart expansion uses this same error tolerance \(\delta_c\),
and all modes of one primary have the same quadratic coefficient.  As in
\cite[Lemma~8.1]{H}, use mode atoms
\[
 \psi_m(y)=\left(
 \zeta^\rho(y)e(c_lp^\rho(y))e(\lambda_my^2+\gamma_my),
 \epsilon_{\rm m}^{1/2}\mathbf e_m\right).
\]
The coefficient sum, mode count, normalized curvature, and rational
heights have logarithms polynomial in the accuracy and localization
parameters \(k,P\).  In particular, for modes of one primary,
\begin{equation}\label{eq:mode-differences-b}
 \gamma_m-\gamma_{m'}=a\cdot(\theta_v,\vartheta),
 \qquad a\in\mathbb Q^{\dim(\theta_v,\vartheta)}
 \text{ of controlled height}.
\end{equation}
Multiplication of the splitting by the fixed integer \(c_l\) changes
only fixed constants in these bounds.

The resulting lift \(\mathcal R_{l,G}\) is one fixed pointwise bounded
operator on augmented inputs.  Indeed, set
\(q_\Phi=\sum_m a_{\Phi,m}(\psi_m)_{\rm phys}
               -(\zeta^\rho)^2\varphi_\Phi\), so that
\(|q_\Phi|\le\delta_c\).  For an augmented input \(F\), its lifted
physical component is
\[
 (\zeta^\rho)^2F_{\rm phys}
       +\sum_{\Phi\ \text{matched}}
             \epsilon_\Phi^{-1/2}F_\Phi q_\Phi,
\]
and its mode coordinate belonging to \((\Phi,m)\) is
\(\epsilon_{\rm m}^{1/2}a_{\Phi,m}
\epsilon_\Phi^{-1/2}F_\Phi\).
Choose \(\delta_c\) small and \(\epsilon_{\rm m}\) small after the
main column budgets.  Cauchy--Schwarz then gives a uniform operator
norm.  The private coordinate of an outer projection output records
its column coefficient, so this operator gives exactly its mode
synthesis.  It is cell-local and independent of the freezing block.

There is also a common periodic multiplier for the four gauges.
For four offsets, put \(d_j=d^{\rho_j}(y_j)\) and
\[
 d_\diamond=\sum_j e_jd_j,\qquad
 k_j=d_j-d_0-b_jd_\diamond.
\]
The \(k_j\) are bounded integer vectors.  Between branch cuts the
lifts are affine and their affine parts cancel in this expression;
hence the carries are locally constant.  The three identities
\(\sum c_jb_j^a=0\), \(0\le a\le2\), give
\begin{equation}\label{eq:gauge-multiplier-b}
 \prod_j e(c_jp^{\rho_j}(y_j))
   =e\left(\sum_jc_j(\theta_zy_j)^TC_zk_j\right).
\end{equation}
For fixed carries the right side is an ordinary character on the
torus with coordinates \((\vartheta y_{l'},\vartheta t)\), for any
anchor \(l'\): \(Q_0\) clears \(C_z\), and changing the anchor uses
integer coefficients \(b_j-b_{l'}\).  Smooth cutoffs slightly larger
than the original ones extend these character pieces by zero at the
cuts.  They give the smooth periodic multiplier, bounded by one,
required in \cite[Section~8.3]{H}.

The number of carries and the derivative bounds acquire only fixed
factors per coordinate or derivative.  Their logarithmic costs remain
polynomial in \(P,k\).  Choose a mode budget \(D_*\ge e^k\),
\(\log D_*\le\operatorname{poly}_b(P,k)\), covering the mode lists,
coefficient sums, heights, inverse private weights, and absolute
Fourier sums, and large enough to dominate the dimensions and all
fixed slope factors.  Only then choose a Fourier cutoff \(B_*\) with
tails at most \(D_*^{-300}\).  This preserves the order of choices
in H; a coefficient-sum budget is not defined through its own
truncation radius.

To control rational divisions at later frequency tests, fix an integer
\begin{equation}\label{eq:Ds-choice}
 D_{\rm s}\ge6,\qquad
 6\mid D_{\rm s},\qquad |b_i-b_j|\mid D_{\rm s}\quad(i\ne j).
\end{equation}
The mode graphs will use shifts in the lattice generated by
\(\vartheta/L_h\), where
\begin{equation}\label{eq:height-parameters-b}
 L_h=D_{\rm s}^{h},\qquad E_h=A_0^{h+1}.
\end{equation}
Here \(E_h\) bounds integer shift indices.  The choice of
\(D_{\rm s}\) is fixed independently of all depth and accuracy
parameters.  From this point, constants indexed by \(b\) may also
depend on this fixed choice of \(D_{\rm s}\).

We recall the scale geometry to make the graph conditions explicit.
In a depth block \([s_0,s_0+K)\), freeze coefficients at
\(u_{s_0}=s_0-4K\), compare on a cell \(J\) at depth
\(v_{s_0}=s_0-2K\), and put \(t_{s_0}=s_0+K\).  A mode has instantaneous frequency
\(\omega_m(x)=2\lambda_mx+\gamma_m\).
On the list frozen at \(u_{s_0}\), join two modes in the PRE graph if
\[
 \sup_{x\in J}|\omega_m(x)-\omega_{m'}(x)
                       +n\cdot\vartheta/L_h|
        \le\Xi/r_{v_{s_0}}
 \quad\text{for some }n\in\mathbb Z^{\dim\theta_z},\quad
 |n|_\infty\le E_h.
\]
The CURRENT graph has the same definition with threshold
\(\Xi/r_{t_{s_0}}\).  Write \(\mathcal E_h(m)\),
\(\mathcal F_h(m)\) for their components.  Both graphs increase
with height, because a witness \(n\) at height \(h\) becomes
\(D_{\rm s}n\) at height \(h+1\), provided \(A_0\ge D_{\rm s}\).
They also increase along successive depth blocks on inherited lists.

Here are parameter choices sufficient for the arithmetic.  Let
\(d_{\max}\ge2\) bound all relevant dimensions, set \(\Delta=10\), and put
\[
 C_*=(10d_{\max}D_*)^{100d_{\max}^4},\qquad
 R'=\lceil\log_2C_*\rceil+2\Delta+10.
\]
Choose the largest height \(J_*\), with \(0\le h\le J_*\), large compared with
\(\delta_*^{-2}d_{\max}R'\log(2D_{\rm m})\), where \(D_{\rm m}\)
is the main-column budget and \(\delta_*>0\) is a sufficiently small
fixed number.  Take \(A_0\) to be a sufficiently large fixed power
of \(C_*D_{\rm s}^{R'}B_*D_*\), and \(\Xi=D_*^{1000}\).
Choose a rational height reserve and a frequency gap with
\[
 \mathcal H>
 (d_{\max}C_*E_{J_*}B_*L_{J_*}Q_0\Xi)^{10},\qquad
 \mathcal T>(d_{\max}\mathcal H)^{Cd_{\max}^4}\Xi.
\]
Finally choose \(K\), polynomial in \(P,k\), after these logarithmic
budgets, so the height-bin packets of \cite[Lemma~3.5]{H} and the
gap \(\mathcal T\) fit inside a block.  Increase it so the pre-lag
Taylor errors are as small as any prescribed fixed inverse power of
\(D_*\).  All degrees can be fixed independently of \(q\).

For a primary with horizontal parameters
\(\Theta=(\theta_v,\vartheta)\), call a block nonexceptional when
every rational vector \(a\) of height at most \(\mathcal H\) satisfies
\begin{equation}\label{eq:tiny-huge-b}
 |a\cdot\Theta|\le\mathcal T^{-1}/r_{v_{s_0}}
 \quad\hbox{or}\quad
 |a\cdot\Theta|\ge\mathcal T/r_{t_{s_0}}.
\end{equation}
The height-bin proof of
\cite[Lemma~8.2]{H} gives \(O(d_{\max})\) exceptional blocks per
primary.  At a nonexceptional block let \(K_{\rm rat}\) be the
rational span of the tiny vectors, and set
\(K_B=\{w:(0,w)\in K_{\rm rat}\}\).
Cramer's rule and the choice of \(\mathcal T\) show that every
height-\(\mathcal H\) vector in \(K_{\rm rat}\) is itself tiny.

\begin{lemma}[Height stabilization]\label{lem:height-b}
For a nonexceptional primary, all but \(O(d_{\max}R')\) heights
\(h\in\{0,\ldots,J_*\}\)
have the following property: two of its modes in one CURRENT
component at height \(h\) are directly adjacent in PRE at height
\(h-\Delta\).
\end{lemma}

\begin{proof}
This is \cite[Lemma~8.3]{H} with the shift lattice
\eqref{eq:height-parameters-b}; we give its divisibility and size
argument.  A simple CURRENT path has at most \(D_*\) edges.
The equal quadratic coefficients of its endpoint modes and
\eqref{eq:mode-differences-b} yield
\[
 a+(0,m_h/L_h)\in K_{\rm rat},\qquad
 m_h\in\mathbb Z^{\dim\theta_z},\quad |m_h|_\infty\le D_*E_h.
\]
Indeed the path bounds its evaluation by \(D_*\Xi/r_{t_{s_0}}\), and
the gap forces the tiny alternative.

Let \(V_i\) be the rational span of the height-\(D_*\) vectors
\(q\) for which \(q+(0,m/L_i)\in K_{\rm rat}\) with
\(|m|_\infty\le D_*E_i\).  Let \(U_i\) be the span of the
integer vectors in \(K_B\) of size at most \(C_*^3D_*E_i\).
These spaces increase.  Discard heights below \(R'\) and those
where either rank changes between \(h-R'\) and \(h\); there are
\(O(d_{\max}R')\) such heights.

At a remaining height, express \(a\in V_{h-R'}\) in a defining
basis \(a=\sum_j\beta_ja_j\), with witnesses \(m_j\).
Cramer's rule bounds the coefficients and their common denominator
by \(C_*\).  Subtracting the relations gives
\[
 m_h-D_{\rm s}^{R'}\sum_j\beta_jm_j\in K_B.
\]
After clearing denominators its size is at most \(C_*^3D_*E_h\),
so it lies in \(U_h=U_{h-R'}=:U\).
The quotient \(\mathbb Z^{\dim\theta_z}/(U\cap\mathbb Z^{\dim\theta_z})\)
is free, since the intersection with the rational space \(U\)
is saturated.  In an integer basis of this quotient,
\[
 [m_h]=D_{\rm s}^{R'}\sum_j\beta_j[m_j].
\]
For each prime \(p\mid D_{\rm s}\), every integral coordinate on
the right has valuation at least
\(R'v_p(D_{\rm s})-\log_2C_*\ge\Delta v_p(D_{\rm s})\).
Thus \([m_h]\) is divisible by \(D_{\rm s}^\Delta\).

Choose an integral lift \(m'\) of this quotient class divided by
\(D_{\rm s}^\Delta\).  Then
\[
 m'-D_{\rm s}^{R'-\Delta}\sum_j\beta_jm_j\in U.
\]
Round its coefficients in an integer spanning set for \(U\)
whose vectors have size at most \(C_*^3D_*E_{h-R'}\).  The adjusted
lift remains integral and satisfies
\[
 |m'|_\infty\le
 D_{\rm s}^{R'-\Delta}\sum_j|\beta_j|\,|m_j|_\infty
   +d_{\max}C_*^3D_*E_{h-R'}\le E_{h-\Delta}.
\]
The last inequality follows from the stated fixed-power choice of
\(A_0\).  No lattice basis or covolume estimate is involved.
Since \(m_h-D_{\rm s}^\Delta m'\in K_B\), the vector
\(a+(0,m'/L_{h-\Delta})\) belongs to \(K_{\rm rat}\) and remains
within the height reserve.  It is tiny by
\eqref{eq:tiny-huge-b}, which gives the required PRE edge.
\end{proof}

We now specify the remaining output imported from
\cite[Section~8]{H}.  Write \(\pi(\mathcal D)\) for the
primary labels represented in a set \(\mathcal D\) of modes.
A mode is good at height \(h\ge\Delta\) when
\[
 \left\lfloor\log_{1.1}|\pi(\mathcal E_{h-\Delta}(m))|\right\rfloor
 =
 \left\lfloor\log_{1.1}|\pi(\mathcal F_h(m))|\right\rfloor
\]
and at least nine tenths of the primaries in \(\mathcal F_h(m)\)
are nonexceptional and satisfy Lemma~\ref{lem:height-b} at \(h\).
Heights below \(\Delta\) are bad.
Call a mode core if it is bad at at most \(\delta_*(J_*+1)\) heights;
call every other mode a linear-side mode.  The increasing graphs and
height stabilization give the coherence of these predicates,
the bad-density count, and hereditary representatives by exactly
\cite[Lemmas~8.4--8.8, Proposition~8.9, and Lemma~8.10]{H}.  Those arguments use component inclusion,
cardinality, projections, and the rational gap; their coordinate
inputs have all been established above.  A minimal PRE or CURRENT
component means a component at height zero.  Good modes in a CURRENT
component have one PRE predecessor at height \(h-\Delta\).  Moreover,
core modes in one minimal CURRENT component have identical good-height
predicates throughout the upper range
\[
 \mathcal H_{\rm up}=
 \{h\in\mathbb Z:\lceil4\delta_*J_*\rceil\le h\le J_*\}.
\]
Each core mode has bad fraction at most \(2\delta_*\) in this range.

We also record the decomposition of the original localized outputs.
Let \(Z_{l,G,s}\) be the current augmented output in role \(l\).
For a basic-\(B\) role in the block beginning at \(s_0\), put
\[
 \widetilde Z_{l,G,s_0}
       =\mathcal R_{l,G} Z_{l,G,u_{s_0}}.
\]
Let \(A_{l,G,s_0}\) and \(V_{l,G,s_0}\) be its exact core and
linear-side mode subtotals, with all private coordinates retained.
Define the full side by
\[
 S_{l,G,s}=Z_{l,G,s}-Z_{l,G,u_{s_0}}+V_{l,G,s_0}
       \qquad(s_0\le s<s_0+K).
\]
In a basic-\(U\) role set \(S_{l,G,s}=Z_{l,G,s}\) and
\(A_{l,G,s_0}=\mathcal E_{l,G,s_0}=0\).
On physical components these definitions give
\[
 Z_{l,G,s}=A_{l,G,s_0}+S_{l,G,s}+\mathcal E_{l,G,s_0}.
\]
For a basic-\(B\) role, the physical error is
\[
 \mathcal E_{l,G,s_0}
   =(Z_{l,G,u_{s_0}})_{\rm phys}
          -(\widetilde Z_{l,G,s_0})_{\rm phys}.
\]

The offsets can be fixed so that every term containing this physical
error is negligible.  The explicit lift, applied to
\(Z_{l,G,u_{s_0}}=\sum_\Phi\beta_\Phi\Phi\), gives
\[
 \mathcal E_{l,G,s_0}
   =(1-(\zeta^\rho)^2)(Z_{l,G,u_{s_0}})_{\rm phys}
                   -\sum_\Phi\beta_\Phi q_\Phi.
\]
The main coefficients and original actual-use cells were fixed before
the offsets.  Their coefficient, size, and path-count bounds, together
with the cutoff expectation, imply the error estimate from
\cite[Section~8, subsection on fixing offsets]{H}:
\[
 \mathbb E_\rho
 \sum_{l,G,I\ {\rm used}}|I|\,
       \|\mathcal E_{l,G,s_0(I)}\|_{2,I}^2
       \le D_{\rm m}^C d_{\max}\delta_c|R|.
\]
With \(\delta_c\) chosen this small in the main budget, fix offsets
satisfying this bound.  All later selections only remove nonnegative
terms; they need not be independent of the offsets.
The linear-side square estimate is uniform in the offsets.
The local absolute \(L^2\) form bound and Cauchy--Schwarz therefore
make the terms containing \(\mathcal E\) negligible, after the same
size stops as in H.  This explains why the subsequent estimates may
use the exact core and full-side arguments.

Here are the quantitative conclusions used below.  A successful
CURRENT component contains a good core mode, whose PRE component
at \(h-\Delta\) is its uniquely determined good predecessor.
For fixed height, role, group, and the pair of floor logarithms of
the predecessor's primary and mode counts, choose a representative
mode inherited under predecessor inclusion.  On comparable actual
use cells \(I,I'\), distinct representatives satisfy H(25), namely
\begin{equation}\label{eq:representative-separation-b}
 |\omega_a(x_*)-\omega_{a'}(x_*)+n\cdot\vartheta/L_h|
 >\frac\Xi2\max(|I|^{-1},|I'|^{-1}),
 \qquad |n|_\infty\le E_h,
\end{equation}
for \(x_*\in I\cap I'\).  On an actual cell, a successful component
\(\mathcal F\) has amplitude
\[
 \sigma_{\mathcal F}^2
  =\sum_{\mathcal D\subseteq\mathcal F}
   \left(\|A_{\mathcal D}\|_{2,I}
       +D_*^{-30}\sum_{\substack{m\in\mathcal D\\m\ {\rm core}}}
                         |\alpha_m|\right)^2 .
\]
Here \(\mathcal D\) runs over minimal CURRENT components,
\(A_{\mathcal D}\) is the full core subtotal in the augmented mode
space, and \(\alpha_m\) is its frozen mode coefficient.
This is H(26); the added coefficient terms absorb uniform Taylor
and Fourier errors.

The Bessel estimates and bad-density count give H(17):
\[
 \sum_{G,s_0}\|V_{l,G,s_0}\|_2^2\le G_0|R|,
 \qquad G_0=\exp\bigl(C_{q,b}(1+\log(2k))^C\bigr).
\]
This is the square hypothesis needed for the side estimates following
Proposition~\ref{prop:short-block-b}.  After the stopping construction
of \cite[Proposition~8.12]{H}, the representatives occurring in a
dyadic amplitude bin \(\sigma_j\) along a retained path satisfy H(28):
\begin{equation}\label{eq:core-path-budgets-b}
 \sum_G|\mathcal A_j^G|
       \le\min(D_*,P^C\sigma_j^{-2}),\qquad
 \sum_{s,G}\|S_{l,G,s}\|_{2,I_s}^2\le P^C.
\end{equation}
Here \(\mathcal A_j^G\) is the set of distinct representative labels
in that bin.  All local
side and component sizes are at most \(P^C\).

For scheduled inputs, these statements are used in the form proved
in \cite[Section~6.3--6.4, especially Lemma~6.3]{V}.  The rows are
cell-local projections composed with the fixed bounded lift
\(\mathcal R_{l,G}\).  Their ordinary Bessel estimates come from
the separated mode construction, and their adjoint main ranges lie
in the usable endpoint spaces established by localization.  V evaluates
the depth input after the final main basic projection.  Its row
composition result therefore applies to the same size, variation,
reading-depth, and endpoint hypotheses.  In particular, the possibly
large mode budget \(D_*\) is not inserted into the strong square
constant \(G_0\).  This is the distinction required by V's side
estimate and by the final choice of \(q\).

\subsection{The side estimate and the near-band gain}

We now estimate the contribution of the full sides to the cell minimum
in \eqref{eq:analytic-base-target}.  The input is the localized expression
\eqref{eq:analytic-localization}, its outgoing-filter identity, the
endpoint variation estimate \eqref{eq:analytic-endpoint-variation},
and the linear-side square bound H(17).  These are precisely the
hypotheses of \cite[Proposition~7.5]{H}; for scheduled inputs their
counterparts are assembled in \cite[Section~6.5]{V}.

Expand each localized argument into its core, full side, and negligible
physical error.  Terms with at least three side factors have a
length-weighted \(p\)-power bound \(P^C|R|\), for any fixed
\(2/3<p<1\).  Indeed Cauchy--Schwarz first sums the group index in
two side roles; writing the three resulting square sizes as
\(a_s,b_s,c_s\), their contribution is bounded using
\[
 \sum_s(a_sb_sc_s)^p
 \le
   \left(\sum_s a_s^2\right)^{p/2}
   \left(\sum_s b_s^2\right)^{p/2}
   \left(\sum_s c_s^2\right)^{p/2}.
\]
This uses \(3p>2\) and the retained path caps; the fourth argument
uses its \(P^C\) size bound.  The group sum is formed before the
power is taken.

In a term with exactly two sides, restore its two core complements
to the whole current arguments \(Z_{c,G,s},Z_{d,G,s}\).
The corrections have a third side.  The outgoing-filter identity
then removes the complementary projections one at a time, transferring
helper averages to one side, say role \(a\).  Its main term is
\[
 H_I(\mathcal D_{a,s}S_{a,G,s},S_{a',G,s},B_c,B_d),
\]
where \(\mathcal D_{a,s}\) is the resulting product of helper
averages.  Restoring the full base pair \(Y_c,Y_d\) introduces
another third-side error.  The pair-smallness from the base reduction
bounds this main term by \(\eta\) times the two side sizes.
Summing costs at most \(\eta P^C|R|\).

At one transfer step, use the outgoing and helper bands defined in
\eqref{eq:analytic-bands}, with these two full sides as its fixed
side operands.  The exact outgoing-minus-helper defect is
\eqref{eq:analytic-band-transfer}.  Its far bands have negligible
error.  For each remaining offset \(h\), \(|h|\le2\), the retained
expression \eqref{eq:analytic-retained-bands} has coefficient
\(-h/P=O(P^{-1})\).  Its paired sum over the band index must remain
inside the cell contribution.

Freeze the complementary band stack and the other complementary
argument at the block's first depth.  Their differences from current
values have square budgets by
\eqref{eq:analytic-endpoint-variation}.  Freeze errors therefore
supply a third strong factor and have the preceding \(p\)-power
bound.  The same argument handles cells where a frozen complement
has large size: its difference from its bounded current value is
then large.  On the remaining ancestor hulls inside a block, both
frozen complements have bounded local \(L^2\) size.

The binary-prefix decomposition in \cite[Lemma~7.4]{H} now writes
the two varying sides, at their uses, as selections from fixed vector
stacks.  The terminal complementary band \(1-Q_{c,P}\) is first
split into the identity and a positive ball term.  Positive outer-ball
ranges have joint group budgets, whereas group-independent identity
and cumulative histories are counted once rather than once per group.
This convention is needed for the following stack energies.
One stack includes the band index.  There are only \(G_0\)
overhead factors.  Individual pad choices stay inside the actual
positive averaged operators: a choice is coupled across depths,
and Jensen's inequality restores the average after estimating it.
They are not additional stack configurations.
The stack energies satisfy H(18):
\[
 \sum_\nu\|\mathbf A_\nu\|_2^2\le PG_0|R|,
 \qquad
 \sum_\nu\|\mathbf B_\nu\|_2^2\le G_0|R|.
\]
Here \(\nu\) records the group, block, and fixed stack configuration.
The functions remain fixed after subsequent cuts.  Independent
Rademacher signs select the two position indices, and a third family
matches the band in \(\mathbf A_\nu\) with the frozen complementary
band stack.  Iterated Khintchine gives the moment hypothesis of
Proposition~\ref{prop:short-block-b}.

For precision, the passage from these energies to a sum of local
stack sizes uses the common hull cuts of \cite[Section~7]{H}.
A hull consists of all ancestors between a use and its maximal use
cell, including their full intervals.  Write \(\Omega_\nu\) for
the union of the original use cells.  The main path and catalog
counts give
\(\sum_\nu|\Omega_\nu|\le M_0|R|\), with
\(\log M_0\le\operatorname{poly}_{q,b}(k)\).
For a family of total stack energy at most \(E|R|\), write
\(\widetilde m_\nu\) for its full hull maxima.  They obey
\[
 \sum_\nu|\{x:\widetilde m_\nu(x)>T\}|
       \le\min(M_0|R|,CE|R|/T).
\]
Stopping at level \(C_1E\) and integrating up to that level costs
at most \(CE(1+\log(2+C_1M_0))|R|\).  Thus only a main-budget
logarithm enters \(G_0\).
For the retained hull define
\[
 m_{A,\nu}(x)=
 \sup_{\substack{I\ {\rm in\ the\ retained\ hull}\\x\in I}}
              \frac1{|I|}\int_I\|\mathbf A_\nu(y)\|^2\,dy,
\]
with value zero for an empty supremum; define \(m_{B,\nu}\) in
the same way.  The weak \(L^1\) estimate for the dyadic maximal
function, followed by truncation and integration, gives common
individual cuts; Markov's inequality for the truncated sum gives
the joint cuts for which these squared hull maxima satisfy
\[
 \sum_\nu m_{A,\nu}\le PG_0,\qquad
 \sum_\nu m_{B,\nu}\le G_0.
\]
The removed roots have arbitrarily small fixed total relative length.
They are first crossings of partial ancestor quantities, so retained
caps hold on whole intervals even in branches containing a later cut.
A sparse decomposition of each retained hull into tree pieces gives
disjoint major subsets of the piece roots.  If \(a_Q,b_Q\) are the
two normalized stack sizes at such a root, then
\[
 \sum_{\nu,Q}|Q|a_Qb_Q
 \le C\int_R\sum_\nu\sqrt{m_{A,\nu}m_{B,\nu}}
 \le P^{1/2}G_0|R|.
\]
Complete each piece by its boundary children.  The two side stacks
then have the required leaf-size bounds, and the frozen complements
have bounded leaf sizes.  Proposition~\ref{prop:short-block-b}
applies to at most \(K\) depths of each piece.  Multiplying by the
transfer coefficient \(O(P^{-1})\) yields the near-band bound
\begin{equation}\label{eq:near-band-b}
 (1+C_N(q))C_{q,b}P^{-1/2}G_0
                    (K+2)^{C_{19}(q-2)}|R|.
\end{equation}
The logarithmic factor of the short-block estimate is included in
\(G_0\).  Arbitrary signs at the original cells are allowed, so this
is an absolute estimate for their coupled band contributions.
No absolute sum over individual paired band indices has been used.

Thus \cite[Proposition~7.5]{H}, with the coefficient-dependent
comparison now proved, gives a complete conditional side bound.
Third-side and freeze errors have \(p\)-power cost \(P^C|R|\);
transferred main terms have absolute cost \(\eta P^C|R|\);
the remaining near bands satisfy \eqref{eq:near-band-b}.
Applying \(\min(\eta,x)\le\eta^{1-p}x^p\) only to the powered
terms gives their contribution \(C_{q,b}\eta^\varepsilon P^C|R|\),
with \(\varepsilon=1-p>0\).  All cuts remove original summands at
first crossings and the data are reconstructed on each normalized
subroot.  This is the root-uniform side estimate used in the closing
argument.  The scheduled version follows from the same fixed stacks
and the row estimates already stated, as in \cite[Section~6.5]{V}.

\subsection{Frequency tests for three or four core arguments}

Fix a retained path point \(x_*\), an actual cell \(I=I_s\) with
\(r=|I|\), a group, and a common successful height
\(h\in\mathcal H_{\rm up}\) for the
core arguments.  Put \(R_s=r^{-1}\) and
\[
 \xi=\vartheta/L_{h-\Delta},\qquad
 \gamma_j^*=\omega_{a_j}(x_*),\qquad B'_h=D_*^5E_{h-\Delta}.
\]
The representative \(a_j\) belongs to the good predecessor of the
component in role \(j\).  Summing a PRE path and using the pre-lag
Taylor bound expresses its selected physical argument, up to the
controlled uniform error, as
\[
 e(c_jp^{\rho_j}(y_j))e(\gamma_j^*(y_j-x_*))
                 g_j(\xi(y_j-x_*)).
\]
The function \(g_j\) is a cutoff times a trigonometric polynomial.
After collecting the common multiplier
\eqref{eq:gauge-multiplier-b}, all Fourier indices in any anchor
coordinates lie in the box of radius \(B'_h\), enlarged by fixed
factors depending on \(b\).  This is H(29) with its complete
coefficient-sum and error bounds.

To compare physical and torus amplitudes, take the horizontal tiny
space \(K_B\) of a nonexceptional primary in a successful component,
and let
\[
 \mathbb H=(K_B\cap\mathbb Z^{\dim\theta_z})^\perp
                   \subseteq\mathbb T^{\dim\theta_z}.
\]
This is a connected torus because the integer lattice being
annihilated is saturated.  The tiny/huge alternative makes physical
interval averages of the relevant characters agree with Haar
averages on \(\mathbb H\), to the prescribed small error.
The component floors above absorb those errors, giving H(30):
\[
 a'_j:=\|g_j\|_{L^2(\mathbb H)}
 \le C_b\left[
   \sum_{\mathcal D\ {\rm selected}}
    \left(\|A_{\mathcal D}\|_{2,I}
       +D_*^{-30}\sum_{\substack{m\in\mathcal D\\m\ {\rm core}}}
                                      |\alpha_m|\right)^2
        \right]^{1/2}
 \le C_b\sigma_{\mathcal F_j}.
\]
The selected sets are whole minimal CURRENT subtotals, by coherence
of the height predicates on \(\mathcal H_{\rm up}\).  For distinct
roles, \(u+(b_j-b_{l'})v\) and \(u+(b_i-b_{l'})v\) have independent
Haar distributions on \(\mathbb H^2\).  Interpolating the six
two-\(L^1\) bounds gives
\begin{equation}\label{eq:core-haar-b}
 \int_{\mathbb H^2}\prod_j|g_j(u+(b_j-b_{l'})v)|\,du\,dv
             \le\prod_j a'_j.
\end{equation}
A bounded progression multiplier changes only the constant.

For four core roles put
\[
 Q=\sum_j\gamma_j^*,\qquad T=\sum_jb_j\gamma_j^*.
\]
Choose, by a fixed rule independent of the mode masks, a reference
\(q'=Q+n\cdot\xi\), \(t'=T+m\cdot\xi\) in the full Fourier
box with \(\max(|q'|,|t'|)\le U_0R_s\), where
\(U_0=D_*^{250}\).  Introduce dyadic parameters by
\begin{equation}\label{eq:four-core-test-b}
 \mu\asymp\min\left(1,\min_j\frac{|t'-b_jq'|}{R_s}\right),
 \qquad
 \nu\asymp\max\left(1,\frac{|q'|}{R_s},\frac{|t'|}{R_s}\right).
\end{equation}
Discard absent references and inner minima below \(U_0^{-1}\).
The four marginals of the normalized kernel make its Fourier
transform vanish on \(t'=b_jq'\).  The mean value theorem near
these lines and rapid Fourier decay give
\(C_{L,b}\mu\nu^{-L}\) for every fixed \(L\).  The tiny/huge
dichotomy puts all moderate frequencies in one pair of character
cosets modulo \(K_B\); summing that pair of cosets is a Haar
integral bounded by \eqref{eq:core-haar-b}.  Thus its contribution
is at most \(C_{L,b}\mu\nu^{-L}\prod_j a'_j\).

For three core roles with side role \(l\), anchor at \(y_l\) and put
\[
 Q=\sum_{j\ne l}\gamma_j^*,\qquad
 T_l=\sum_{j\ne l}(b_j-b_l)\gamma_j^*.
\]
Choose a moderate reference \(T_l+m\cdot\xi\) and write
\(\beta=r(T_l+m\cdot\xi)\).  Set
\(\mu\asymp\min(1,|\beta|)\),
\(\nu\asymp\max(1,|\beta|)\), with the same small and large
discards.  Partial Fourier integration of the kernel in the step
variable gives a function \(W_\beta\) whose fixed-order derivatives
are bounded by \(C_{L,b}\mu\nu^{-L}\).  At zero this uses the
\(l\)-marginal; at large frequency it uses integration by parts.
Duality in \eqref{eq:core-haar-b} bounds the remaining horizontal
coefficient sequence in \(\ell^2\).

More explicitly, average a fresh offset cutoff
\(\zeta^{\rho_l}/c_\zeta\) in the side, where
\(c_\zeta=\mathbb E_{\rho_l}\zeta^{\rho_l}\ge1/2\) is constant.
This insertion is an exact identity in expectation.
Keep the bounded factor \(c_\zeta^{-1}\) outside the estimate and set
\(S'=\zeta^{\rho_l}e(-c_lp^{\rho_l})S_{\rm phys}\) on \(I\).
Its normalized \(L^2\) norm is at most that of the side.  Write
\(\widehat S'_I(k')\) for its normalized Fourier coefficients.
The horizontal pairings decay in
\(|k'+r(Q+n\cdot\xi)|\).  Schur's estimate on a dyadic shell
of radius \(\Lambda\), followed by averaging the fresh offset,
gives H(32) with
\begin{equation}\label{eq:side-windows-b}
 b_\Lambda^2=
 \mathbb E_{\rho_l}
 \sum_{\substack{k'\in\mathbb Z:\\
 |k'+r(Q+n\cdot\xi)|\le2\Lambda\ \text{for some}\ n,\\
 |n|_\infty\le B'_h}}
             |\widehat S'_I(k')|^2.
\end{equation}
The contribution is bounded by
\(C_{L,b}\mu\nu^{-L}\prod_{j\ne l}a'_j
\sum_{1\le\Lambda\le U_0}\Lambda^{-L}b_\Lambda\),
where the sum is dyadic.  Notice that the windows use the whole
Fourier box and are independent of the mode masks.

The error estimates in \cite[Section~9.1--9.2]{H} now apply with
these explicit tests.  They use the absolute local form estimate,
Fourier tails, the pre-lag, and the rational gap, all with reserves
already chosen above.  Per tuple the error is bounded by a fixed
multiple of \(D_*^{-100}\) times the absolute core coefficient
sums, and by the side size if a side is present.  The common-height
argument samples heights independently and uniformly from
\(\mathcal H_{\rm up}\) until all roles are good; its first-success
weights sum to one on each pure tuple.  A core mode has bad fraction
at most \(2\delta_*\), so the sum over failure assignments converges.
Consequently these errors sum absolutely to \(O(e^{-k}|R|)\), and
the principal cell terms have the form H(33):
\begin{equation}\label{eq:tested-cell-b}
 C_{L,b}\mu\nu^{-L}\Lambda^{-L}
                       \left(\prod_{j=0}^3\sigma_j\right)N_s^G.
\end{equation}
Here the \(\sigma_j\) are dyadic component-amplitude bins, with
the side bin belonging to \(b_\Lambda\); \(N_s^G\) counts the
representative tuples passing the tests.  For four cores,
\(\Lambda=1\).  We next bound these counts across scales.

\subsection{Sparse mass for separated tests}

The finite counting statement is useful independently of the mode
construction.  Its integer divisibility is the reason for the
choice \eqref{eq:Ds-choice}.

\begin{lemma}[Sparse mass]\label{lem:sparse-b}
Let \(J\subseteq\{0,1,2,3\}\) consist of three or four core roles.
In each role let a finite label set \(\mathcal A_i\) be identified
injectively with real numbers, with \(|\mathcal A_i|\le M\),
\(M\ge1\).  Let \(d\ge0\) be an integer.  Fix
\(\zeta\in\mathbb R^d\), \(B,L\ge1\), \(Q,\Sigma,R_0>0\), and
\(0<\mu\le1\le\nu\).  Every shift vector below belongs to
\(\mathbb Z^d\).  For finitely many distinct integers \(s\), at
dyadic scales \(R_s=R_02^s\), let
\(\mathcal E_s\subseteq\prod_{i\in J}\mathcal A_i\)
and \(N_s=|\mathcal E_s|\).  Assume the following five conditions.
\begin{enumerate}
\item Two labels in different core roles determine at most one full
tuple, and that pair occurs at at most \(L\) scales.
\item If there is a side role, there are \(w_s\ge0\),
\(\sum_s w_s\le M\), such that every fixed core label is incident
to at most \(w_s\) tuples at scale \(s\).
\item Distinct same-role labels occurring at scales \(s,t\) satisfy
\[
 |a-a'+n\cdot\zeta|>\Sigma\max(R_s,R_t)
       \quad(|n|_\infty\le Q).
\]
\item For every triple \(T\subseteq J\), with excluded role
\(l'\in\{0,1,2,3\}\setminus T\) and
\(\alpha_i=b_i-b_{l'}\), each tuple has a witness
\(|n|_\infty\le B\) such that
\[
 \mu R_s\le\left|\sum_{i\in T}\alpha_i a_i
                       +D_{\rm s}^2 n\cdot\zeta\right|\le\nu R_s.
\]
\item At every occupied scale,
\( |n\cdot\zeta|\notin[\mu R_s/1000,1000\nu R_s]\)
for all \(|n|_\infty\le Q\).
\end{enumerate}
There are constants \(C_b,C_{{\rm res},b}\) such that, if
\[
 Q\ge C_{{\rm res},b}(d+1)^2B,\qquad
 \Sigma\ge C_{{\rm res},b}(d+1)^2\nu,
\]
then, with \(\ell=1+\log_2(2+(d+1)^2\nu/\mu)\),
\[
 \sum_{s:N_s\le\delta M^2}N_s
      \le C_b\ell L^{29}\delta^{1/2080}M^2
       \qquad(0<\delta\le1).
\]
The exponents and constants are independent of the dimension and
the number of labels and scales, apart from the indicated factors.
\end{lemma}

\begin{proof}
We follow the finite argument of \cite[Lemma~9.1]{H}, spelling out
the arithmetic and the dimension dependence.  The proof first covers
most label differences by neighborhoods associated with later scales.
Three such future covers then place sums of triples in a thin slab.
Separated anchor labels reduce that slab's cover from three factors
of \(M\) to two, while the pure-shift gap makes many translated
triples distinct inside the cover.  This comparison forces small
mass at sparsely occupied scales.

Approximate all finite
real data on a sufficiently fine common grid, taking approximants
on a fixed divisible multiple of its step so divisions by the
\(\alpha_i\) remain on the grid.  For example, a multiple of
\(D_{\rm s}^2\) suffices.  Choosing the grid after the data preserves
separation and annular inequalities with fixed relaxed margins and
leaves a narrower pure-shift gap.  Round scalar widths upward by at
most one grid step.  The estimates below are uniform in that step.

Restrict to \(N_s\le\delta M^2\).  Pair determination gives total
mass at most \(LM^2\).  For parameters \(\delta\le\alpha/2\)
and \(0<\kappa<1\), greedily partition the scales in decreasing
order of \(R_s\) into counting blocks of mass at most
\(\alpha M^2\).  There are at most \(J_0=3L/\alpha\) blocks.
For a block \(\mathfrak b\), let \(R_{\mathfrak b}^+\) be its
largest scale, take \(N_0=\lceil C_bB\rceil\) with
\(C_b>32D_{\rm s}\), and choose a grid width
\(\epsilon_{\mathfrak b}\asymp\nu R_{\mathfrak b}^+\).
Use the finite neighborhoods
\[
 U_{\mathfrak b}(t)=
 \{D_{\rm s}n\cdot\zeta+u:
 |n|_\infty\le tN_0,\ u\text{ on the grid},\
 |u|\le t\epsilon_{\mathfrak b}\}.
\]
They are symmetric and satisfy
\(U(t)+U(t')\subseteq U(t+t')\).

Fix a role and a counting block with uncovered mass exceeding
\(\kappa M^2\).  Project onto this role and another core role,
retaining one scale witness per pair, and color edges by a third
core label.  The graph has at least \((\kappa/L)M^2\) edges and
a proper coloring with at most \(M\) colors.  The elementary
four-step-walk observation in the proof of \cite[Lemma~9.1]{H}
gives a subset \(B_0\) of at least \(c\eta M\) vertices,
\(\eta=\kappa/L\), such that any two have at least
\(c\eta^5M^3\) four-step walks with distinct ordered color tuples.
Repeated vertices are allowed.
That observation is purely a graph statement: choose a random
neighborhood, remove vertices with many low-codegree partners, and
count the two choices of a common neighbor on either side of a
retained midpoint.

Alternating the four annular relations along a walk cancels the
intermediate labels.  Dividing by the nonzero integer \(\alpha_i\)
gives
\[
 a-a'\in q(c_1,c_2,c_3,c_4)+U_{\mathfrak b}(1).
\]
Here \(q\) depends only on the ordered four edge colors; these
\(c_i\) are graph colors, unrelated to the fixed quadratic weights.
The scalar error after division is at most
\(4\nu R_{\mathfrak b}^+\), and
\(D_{\rm s}^2/\alpha_i\) is an integer multiple of
\(D_{\rm s}\), with its index covered by \(N_0\).
Packed differences outside \(U_{\mathfrak b}(2)\) cannot use the
same color tuple.  The \(M^4\) possible color tuples therefore
bound their number by \(C\eta^{-5}M\).  Repeating after removing
\(B_0\) covers all but \(\kappa M^2\) occurrences in each role
and block by at most
\(K_1=C(L/\kappa)^6\) subsets, each with its difference set
covered by \(K_1M\) translates of \(U_{\mathfrak b}(2)\).

Discard uncovered occurrences.  Call a core label future when it
belongs to a cover subset in a later counting block.  Nonfuture
label sets are disjoint over blocks.  Let \(m_i^{\mathfrak b}\)
be their cardinalities in core role \(i\).  Treat the side as
nonfuture, with \(m_l^{\mathfrak b}=\sum_{s\in\mathfrak b}w_s\).
Occurrences with fewer than three future roles have two nonfuture
roles.  Their total mass is at most \(C\sqrt{L\alpha}M^2\), by
summing
\(\min(\alpha M^2,Lm_i^{\mathfrak b}m_j^{\mathfrak b})\)
and Cauchy--Schwarz.  The discarded uncovered mass is at most
\(4J_0\kappa M^2\).

Assign each remaining occurrence three future roles and one future
cover subset in each.  A class fixes its current counting block
and its three future block/subset choices.  There are at most
\(4J_0^4K_1^3\) classes.  Fix one of mass \(\rho M^2\), with
future subsets \(B_1,B_2,B_3\).
The subscripts now enumerate those roles; their coefficients remain
the corresponding \(\alpha_i\).  Set
\(\Gamma=C_{{\rm gap},b}(d+1)^2\nu/\mu\).
If \(\rho\le C_b\delta\ell\), retain this bound.  Otherwise split
the mass into four consecutive ranges by quartiles, omitting the
at most three scale bins crossing a cut.  Each range retains at
least \(\rho M^2/8\) mass.  The first provides anchor sets
\(B_{i0}\subset B_i\) of size at least \(c\rho M/L\).
Let \(R^\#\) be the largest scale of the second range; these
anchors are separated modulo reserved shifts at scale \(R^\#\).
The fourth range gives a set \(F\) of at least \(c\rho M^2/L\)
distinct triples projected onto these three roles.
The third range supplies the scale margin, and a greedy thinning
of the second gives \(q\ge c_b\rho/(\delta\ell)\) triples
\(t_v\), also projected onto these roles, at scales separated by
factors at least \(\Gamma\).
All future-block scales and fourth-range scales are smaller than
these by at least the required factor \(\Gamma\).

Write \(U_i\) for the neighborhoods of the three future subsets.
Subdividing index and scalar intervals covers \(U_i(t)\) by at
most \((2t+1)^{d+1}\) translates of \(U_i(1)\).  Hence there is
an odd \(H_i\le100(d+1)^2\) with
\(|U_i(H_i+2)|\le2|U_i(H_i)|\).  The preliminary exponential
cover is used only to find this small-growth index, and is never
paid in the cardinality estimate.  The translates
\(a+U_i(H_i)\), \(a\in B_{i0}\), are disjoint: an intersection
would contradict the anchor separation with a shift index
\(O_b((d+1)^2B)\) and scalar error
\(O_b((d+1)^2\nu R^\#)\).

Set \(Z_i=B_{i0}+U_i(H_i)\).  The difference covers give
\( |Z_i-B_i|\le C(K_1L/\rho)|Z_i|\).
Choose a nonempty \(X_i\subset Z_i\) minimizing
\(r_i=|X_i-B_i|/|X_i|\).  Petridis's minimal-growth subset
argument, in the form reproduced in \cite[Lemma~9.1]{H}, gives
\( |X_i-mB_i|\le r_i^m|X_i|\).
One anchor translate contains a subset \(Y_i\subseteq X_i\)
of at least \(|X_i|/M\) points.  If \(\mathcal P\subseteq mB_i\)
has all distinct differences outside \(U_i(2H_i)\), the sets
\(Y_i-p\), \(p\in\mathcal P\), are disjoint and lie in
\(X_i-mB_i\).  Thus \(|\mathcal P|\le r_i^mM\).
A maximal packing covers \(mB_i\) by that
many translates of \(U_i(2H_i)\).  Only \(m=2,3\) is used.

The outputs \(t_v+F\) lie in the slab
\[
 (x_1,x_2,x_3)\in2B_1\times2B_2\times2B_3,\qquad
 \left|\sum_i\alpha_ix_i+D_{\rm s}^2n\cdot\zeta\right|
                  \le2\nu R^\#,\quad |n|_\infty\le2B.
\]
Put \(r_0=CK_1L/\rho\).  Cover the first two sumsets by
\(r_0^2M\) boxes each.  Above a fixed pair, dividing the slab
equation by \(\alpha_3\) places the third coordinate in
\[
 c+\{n\cdot\zeta:|n|_\infty\le C_b(d+1)^2B\}
       +[-C_b(d+1)^2\nu R^\#,C_b(d+1)^2\nu R^\#].
\]
Every shift remains integral: both \(D_{\rm s}^2/\alpha_3\)
and \(D_{\rm s}\alpha_i/\alpha_3\) are integers.
If \(\mathcal P\) is a packing of the local third-coordinate
set outside \(U_3(2H_3)\), then
\(B_{30}+\mathcal P\subset3B_3\) is still a packing.
Collisions with different anchors are excluded by the same shift
and scalar bounds.  Thus
\( |B_{30}|\,|\mathcal P|\le r_0^3M\).
The entire slab is covered by at most
\(CK_1^7(L/\rho)^8M^2\) product boxes.

Each such box contains at most one of the indexed outputs
\(t_v+F\).  Within one translation this follows from label
separation.  Across translations, let \(R_{\rm big}\) be the
larger selected scale.  Its annular residual has magnitude at
least \(\mu R_{\rm big}\); the other three residuals together
are at most \(\mu R_{\rm big}/10\).  Membership in one box
adds scalar error at most
\(C_b(d+1)^2\nu\max_i R_{{\rm future},i}^+
\le\mu R_{\rm big}/10\), by the choice of \(C_{{\rm gap},b}\).
The resulting pure shift has magnitude between
\(4\mu R_{\rm big}/5\) and \(3\nu R_{\rm big}\), with index
\(O_b((d+1)^2B)\le Q\), contradicting the forbidden gap.

Comparing \(q|F|\) with the slab cover yields
\(\rho^{10}\le C_b\delta\ell K_1^7L^9\), and hence
\(\rho\le C_b\ell K_1L\delta^{1/10}\).
Including all choices and losses gives
\[
 M^{-2}\sum_{s:N_s\le\delta M^2}N_s
 \le C_b\left(\frac{L\kappa}{\alpha}
       +\sqrt{L\alpha}
       +\ell L^{29}\alpha^{-4}\kappa^{-24}\delta^{1/10}\right).
\]
Take \(\alpha=\delta^{1/1040}\), \(\kappa=\alpha^2\) for
sufficiently small \(\delta\), and use the crude mass bound for
the remaining values.  This proves the stated exponent \(1/2080\).
\end{proof}

\subsection{Application of the count and the resulting core bound}

We verify the hypotheses rather than apply the original arithmetic
form of \cite[Proposition~9.2]{H}.  Fix a path, height, role types,
frequency bands, and amplitude bins in \eqref{eq:tested-cell-b}.
Identify each representative label with its value
\(\omega_a(x_*)\), which is injective in a fixed role and type by
\eqref{eq:representative-separation-b}.  Define
\[
 M_i^G=|\mathcal A_i^G|\quad\text{for core roles},\qquad
 M_l^G=\sigma_l^{-2}\sum_s\|S_{l,G,s}\|_{2,I_s}^2
          \quad\text{for a side},\qquad M'=\max_iM_i^G.
\]
Groups with a nonzero count have \(M'\ge1\).

Two fixed core labels determine the others.  For four cores,
subtract the equations for \(Q,T\) of two candidate tuples.
The two unknown label differences solve a system with determinant
\(b_j-b_i\), modulo bounded rational shifts and errors
\(O_b(\nu\max(R_s,R_t))\).  For three cores the equation for
\(T_l\) determines the last label, dividing by \(b_j-b_l\).
The shift reserve and separation force these differences to vanish.
For every core triple with excluded role \(l'\), the tests also
give the annulus
\begin{equation}\label{eq:annulus-application-b}
 c_b\mu R_s\le
 \left|\sum_i(b_i-b_{l'})\omega_{a_i}(x_*)+n\cdot\xi\right|
 \le C_b\nu R_s,\qquad |n|_\infty\le C_bB'_h.
\end{equation}
For four cores this is \(t'-b_{l'}q'\); for three it is the
\(T_l\) test.  Once the labels are fixed, the height-bin lemma
\cite[Lemma~3.5]{H} bounds their scale occurrences by
\(L_0\le P^C(1+\log(\nu/\mu))^C\).

If one role is a side, windows \eqref{eq:side-windows-b} for
different triples incident to a fixed core label are disjoint.
An overlap gives both a relation for their \(Q\) values with
error \(O(\Lambda R_s)\) and their two \(T_l\) relations.
Solving the resulting two-by-two system again divides by a
\(b_j-b_i\); separation then identifies the tuples.  Couple the
fresh offset across these windows.  Parseval shows their incidence
count is at most
\(w_s=\sigma_l^{-2}\|S_{l,G,s}\|_{2,I_s}^2\), giving the
side-budget hypothesis.

Here are the promised reserve checks for all these divisions.  In
Lemma~\ref{lem:sparse-b} take
\[
 \zeta=\xi/D_{\rm s}^2
       =D_{\rm s}^{\Delta-2}\vartheta/L_h,\qquad
 Q_{\rm res}=\left\lfloor E_h/D_{\rm s}^{\Delta-2}\right\rfloor,
 \qquad \Sigma=\Xi/2.
\]
Every divided \(\xi\)-shift is integral in this finer lattice,
because each determinant divides \(D_{\rm s}\).
The fixed-power choice of \(A_0\) ensures
\[
 Q_{\rm res}\ge
 C_{{\rm res},b}(d+1)^2C_bB'_h,
 \qquad d=\dim\theta_z.
\]
Also \(\nu,\Lambda\le C_bD_*^{250}\), while
\(\Xi=D_*^{1000}\), so the separation reserve holds.  All
rational shifts remain below \(\mathcal H\).
Every successful component contains a nonexceptional primary,
even if its representative is exceptional.  Applying
\eqref{eq:tiny-huge-b} to its pure horizontal shifts excludes the
entire intermediate interval required by the pure-shift hypothesis,
since \(\mu\ge U_0^{-1}\) and the interval length \(|I_s|\) lies
between \(r_{t_{s_0}}\) and \(r_{v_{s_0}}\).  This checks all five hypotheses.

Lemma~\ref{lem:sparse-b}, pair determination, and the side incidence
bound therefore give the two counts H(34): for some fixed
\(p_0<1\),
\begin{align}
 \sum_sN_s^G&\le L_0\min_{i\ne j}M_i^GM_j^G,\label{eq:count-one-b}\\
 \sum_s(N_s^G)^{p_0}
 &\le P^C(1+\nu+\mu^{-1}+\Lambda)^C(M')^{2p_0}.
 \label{eq:count-power-b}
\end{align}
Indeed the sparse-mass exponent \(c'=1/2080\) allows any
\(1-p_0<c'\).  On the bins
\(2^{-j-1}(M')^2<N_s^G\le2^{-j}(M')^2\), apply the sparse
bound and sum the geometric series with ratio
\(2^{1-p_0-c'}<1\).

Finally the amplitude and group summation of
\cite[Proposition~9.3]{H} uses precisely
\eqref{eq:tested-cell-b}, \eqref{eq:core-path-budgets-b}, and
\eqref{eq:count-one-b}--\eqref{eq:count-power-b}.  Interpolate the
two counts with a sufficiently small fixed weight \(\theta'>0\),
giving power \(p=(1-\theta')+\theta'p_0\).  The path square
budgets sum the group factors, and the core count cap \(D_*\)
controls the common amplitude shift.  The three amplitude-gap sums
are geometric, as in that proposition.  The frequency losses after
interpolation are at most
\(\mu^{-C\theta'}\nu^{C\theta'}\Lambda^{C\theta'}\);
choose \(\theta'\) small and then the derivative order \(L\)
large so they are summable against
\(\mu^p\nu^{-Lp}\Lambda^{-Lp}\).
Applying the original cell minimum before subadditivity yields
\begin{equation}\label{eq:core-final-b}
 C_{q,b}\eta^{\varepsilon}P^{C_4}|R|+O(e^{-k}|R|)
\end{equation}
for the terms with at least three core arguments, with fixed
\(\varepsilon>0\) and \(C_4\) independent of \(q\).
The same calculation is the core step of \cite[Section~6.6]{V}:
its frozen scheduled inputs provide exactly the amplitudes,
representatives, and same-cell side incidence verified above.
Together with Proposition~\ref{prop:short-block-b} and the side
square bound, this supplies the side and core estimates required
to close the proof of Theorem~\ref{thm:local}.
